\section{命题与逻辑连接词}
\label{sec:01-Mathematical-Language/01-Logic/01}

一份模型报告里写着：``新版本的效果更好。''

这句话读起来毫无问题，却没办法进入任何推理链条。更好是指哪个指标？在哪批数据上？提升多少才算好？换个人看同一组结果，完全可以给出相反的结论。它不是错的，它是无从判断对错的。

对照一句明摆着的错话：``整数 15 是偶数。''这句话是假的，但它合格——任何人拿到它，都会给出同一个判断。

数学要求的第一件事就是这个。一句话必须先具备确定的真假，才谈得上对不对。做不到这一点的句子，无论听起来多有道理，都不能当作推理的起点。这不是吹毛求疵：后面所有的定理、所有的证明，都建立在``每一步的真假不因人而异''之上，一旦某一步含混，整条链子就断了。

介于两者之间的是 \(x>3\) 这样的句子。它自己不真不假，要等 \(x\) 拿到具体值，或者说清楚 \(x\) 在哪个范围里取，才能判断。这种带自由变量的句子叫谓词。谓词不是残次品——数学里绝大多数有价值的断言，都是谓词配上量词组装出来的，第三篇会专门处理这件事。

\subsection{连接词：把判断拼装成更大的判断}

设 \(P,Q\) 是两个命题。五个基本连接词分别是 \(\neg P\)（非 \(P\)）、\(P\land Q\)（\(P\) 且 \(Q\)）、\(P\lor Q\)（\(P\) 或 \(Q\)）、\(P\to Q\)（若 \(P\) 则 \(Q\)）、\(P\leftrightarrow Q\)（\(P\) 当且仅当 \(Q\)）。它们的取值只依赖 \(P,Q\) 的真假，不依赖 \(P,Q\) 讲的是什么内容——这正是它们能被机械处理的原因。

真正容易出事的是``或''。点餐时说``咖啡或茶''，潜台词是二选一；数学里的 \(P\lor Q\) 默认允许两个都成立，只要求至少一个成立。想表达``恰好一个''，必须专门写异或。

否定怎样穿过``且''与``或''，由 De Morgan 律规定：

\[
\neg(P\land Q)\equiv(\neg P)\lor(\neg Q),
\qquad
\neg(P\lor Q)\equiv(\neg P)\land(\neg Q).
\]

这两条经常在写筛选条件时被违反。假设要把``年龄不小于 18 且消费不低于 100 元''的用户排除掉，剩下的人满足的条件是``年龄小于 18 \emph{或} 消费低于 100 元''。写成``年龄小于 18 且消费低于 100 元''，就会把只违反一项的人留在集合里。否定一个合取，得到的是析取，反过来也一样：中间那个词必须跟着翻转。

\subsection{蕴含只有一种失败方式}

在 \(P\to Q\) 里，\(P\) 叫前件，\(Q\) 叫后件。这两个名字只描述它们在箭头的哪一边，既不表示 \(P\) 发生得更早，也不表示 \(P\) 是 \(Q\) 的原因。

把 \(P\to Q\) 读成一个承诺：只要 \(P\) 成立，\(Q\) 就必须跟着成立。承诺被打破，当且仅当有人拿出一个 \(P\) 成立而 \(Q\) 不成立的实例。其余三种情形都没有构成违约。

\begin{center}
\begin{tikzpicture}[x=2.5cm,y=1.35cm,font=\small]
  \fill[black!5]  (0,1) rectangle (1,2);
  \fill[red!14]   (1,1) rectangle (2,2);
  \fill[black!5]  (0,0) rectangle (1,1);
  \fill[black!5]  (1,0) rectangle (2,1);
  \draw[black!45] (0,0) rectangle (2,2);
  \draw[black!45] (1,0) -- (1,2);
  \draw[black!45] (0,1) -- (2,1);
  \node[align=center] at (0.5,1.5) {承诺兑现};
  \node[align=center] at (1.5,1.5) {\textbf{唯一的反例}};
  \node[align=center,black!55] at (0.5,0.5) {前提未触发};
  \node[align=center,black!55] at (1.5,0.5) {前提未触发};
  \node[above] at (0.5,2) {\(Q\) 真};
  \node[above] at (1.5,2) {\(Q\) 假};
  \node[left]  at (0,1.5) {\(P\) 真};
  \node[left]  at (0,0.5) {\(P\) 假};
\end{tikzpicture}

\emph{四种真假组合里，只有右上角那一格能推翻 \(P\to Q\)；下面两格连前提都没满足，没有资格出庭作证。}
\end{center}

下面两格常让人别扭：前提都假了，凭什么整个蕴含还算真？换个角度就顺了。这条命题禁止的是一类对象——满足前提却不满足结论的对象。一个连前提都不满足的东西，根本进不了被告席。

拿``若整数 \(n\) 能被 4 整除，则 \(n\) 是偶数''来说，整数 3 既不被 4 整除也不是偶数，但它不是反例。合格的反例必须是``4 的倍数却不是偶数''的数。这种数不存在，所以命题成立。

同一件事也可以写成等价式 \(P\to Q\equiv\neg P\lor Q\)：要么前提落空，要么结论成立，二者有其一，承诺就没被打破。

\subsection{等价是一张改写许可证}

两个复合命题，如果在基本命题的每一种真假组合下结果都相同，就称它们逻辑等价。注意这个要求有多强：不是``通常一起对''，而是一种例外都不允许。真值表就是逐行核对这件事的工具，变量少的时候极好用。

验证过的等价式随后可以当改写规则使用。把 \(P\to Q\equiv\neg P\lor Q\) 和 De Morgan 律接起来：

\[
\neg(P\to Q)\equiv\neg(\neg P\lor Q)\equiv P\land\neg Q.
\]

这个结果值得记住，因为它直接给出了反例的形状：要反驳一条蕴含，就去找一个使前件成立、后件失败的对象——不多不少刚好这一种。后面读到``某某定理不成立''的讨论时，找的都是这个东西。

试一个稍复杂的：``若 \(n\) 是 6 的倍数，则 \(n\) 既是 2 的倍数又是 3 的倍数。''记 \(P(n):6\mid n\)、\(Q(n):2\mid n\)、\(R(n):3\mid n\)，命题形如 \(P(n)\to(Q(n)\land R(n))\)，其否定是

\[
P(n)\land\bigl(\neg Q(n)\lor\neg R(n)\bigr).
\]

于是反例必须是 6 的倍数，同时至少漏掉 2 和 3 中的一个。写出 \(n=6k=2(3k)=3(2k)\) 就知道这种数不存在：两个结论是同一个因式分解的两种读法。请注意这里发生了什么——把否定推到底之后，证明该往哪走已经自动显形了。

真值表也有它够不着的地方。它把 \(P,Q\) 当成不可再分的真假值，一旦命题内部涉及整数、函数或无限集合，逐行枚举就失效了，必须转而研究对象本身的结构。这条分界线是命题逻辑与谓词逻辑的分界线。

\subsection{三种反复出现的误用}

\begin{itemize}
\tightlist
\item
  把 \(P\to Q\) 当成 \(P\) 导致 \(Q\)。逻辑蕴含只约束真假的搭配，不承诺任何机制。冰淇淋销量高时溺水人数也高，蕴含关系成立，因果关系不成立。
\item
  由 \(P\to Q\) 与 \(Q\) 推出 \(P\)（肯定后件）。下雨地面会湿，但地面湿也可能是洒水车。模型在测试集上表现好，不等于它学到了正确的规律。
\item
  由 \(P\to Q\) 与 \(\neg P\) 推出 \(\neg Q\)（否定前件）。前提没触发，结论仍然可以由别的途径成立。
\end{itemize}

后两条之所以顽固，是因为它们在日常对话里往往能猜对——说话人通常还暗示了反方向。数学里没有这层暗示，方向必须写明。

单向箭头因此有了两种读法：把 \(P\) 看成保证 \(Q\) 的手段，还是看成 \(Q\) 得以成立的前提？下一篇把这条箭头拆成充分条件与必要条件，并说明为什么``仅当''指向的方向和多数人的第一直觉相反。
